% ECONOMICS_PROBLEM: {"id": "C73-17", "title": "Maxmin equality with the concavified nonrevealing value", "jel_code": "C73", "source_id": "OP-0168"}
% !TeX program = xelatex
\documentclass[11pt,a4paper]{article}
\usepackage[margin=23mm]{geometry}
\usepackage{fontspec}
\setmainfont{Latin Modern Roman}
\usepackage{amsmath,amssymb,mathtools}
\usepackage{xcolor,enumitem,booktabs,longtable,array,xurl,hyperref}
\definecolor{muted}{HTML}{53636C}
\hypersetup{colorlinks=true,urlcolor=blue,linkcolor=blue}
\setlength{\parindent}{0pt}
\setlength{\parskip}{5pt}
\newcommand{\R}{\mathbb R}
\newcommand{\E}{\mathbb E}
\newcommand{\N}{\mathbb N}
\newcommand{\ind}{\mathbf 1}
\newcommand{\Var}{\operatorname{Var}}
\newcommand{\Cov}{\operatorname{Cov}}
\newcommand{\supp}{\operatorname{supp}}
\DeclareMathOperator*{\argmin}{arg\,min}
\DeclareMathOperator*{\argmax}{arg\,max}
\newcommand{\entrymeta}[3]{{\sffamily\small\color{muted}\textbf{\texttt{#1}}\quad|\quad #2\par #3\par}}
\newcommand{\Problemref}[1]{\texttt{#1}}
\newcommand{\JELref}[2]{\texttt{#2} (JEL \texttt{#1})}
\begin{document}
\section*{Maxmin equality with the concavified nonrevealing value}
\textbf{Problem ID:} \texttt{C73-17}\quad \textbf{JEL:} \texttt{C73}\par
% BEGIN ECONOMICS STATEMENT
\label{problem:OP-0168}
% GAME_THEORY_PROBLEM: {"local_id": "GT-038", "original_number": 38, "kind": "P", "entry_kind": "statement_only", "published": false, "review_required": true, "category": "game_theory"}
\label{game:GT-038}
{\sffamily\small\color{muted}
{\sffamily\small\color{muted}\textbf{GT-038} | Repeated games with incomplete information\par P | Source-reported characterization question; a tail-measurable subclass is known\par}
\par}
\subsection*{Definitions and mathematical statement}
Let $K,I,J$ be finite. Nature draws $k\sim p\in\Delta(K)$ and reveals it only to the maximizing player; the players observe every subsequent joint action. For bounded Borel $f:K\times(I\times J)^{\mathbb N}\to\mathbb R$, let $\underline v_f(p)$ be the supremum over informed behavioral strategies $\sigma(k,h)$ of the infimum over uninformed strategies $\tau(h)$ of the expected payoff $f$.

If the maximizing strategy also depends only on the public action history, it is nonrevealing. Define the complete-information functional
\[
 f_p(\omega)=\sum_{k\in K}p(k)f(k,\omega).
\]
The associated bounded-Borel Blackwell game has a value, denoted $u_f(p)$; it is the nonrevealing value. Define its concave envelope by
\[
 (\operatorname{cav}u_f)(p)=
 \sup\left\{\sum_{r=1}^{R}\lambda_r u_f(q_r):
 R\geq1,\ q_r\in\Delta(K),\ \lambda_r\geq0,\ 
 \sum_r\lambda_r=1,\ \sum_r\lambda_rq_r=p\right\}.
\]
For fixed alphabets, characterize bounded Borel $f$ satisfying
\[
 \forall p\in\Delta(K),\qquad
 \underline v_f(p)=(\operatorname{cav}u_f)(p).
\]
The desired conditions should refer to the functional's structure and not simply repeat the equality to be characterized.
\subsection*{Short English statement}
Exactly when is the informed player's guaranteed payoff the concavification of the nonrevealing value?
\subsection*{Variants and scope boundaries}
Nonrevealing strategies cannot condition directly on $k$. The equality concerns the maxmin, not necessarily the minimax or existence of a game value. Dependence on an initial finite block can matter for general Borel functionals. A one-prior characterization and an effective recognition algorithm would be different targets.
\subsection*{Source relationship and status}
[S1], Theorem 2.5, establishes this equality for its tail-measurable class. Section 5 asks for its characterization beyond that class. The known tail-measurable theorem is therefore a boundary of the open target, not itself an open conjecture.
\subsection*{References}
\begingroup\footnotesize\raggedright
\textbf{[S1]} Gil Bar Castellon Koltun, Ehud Lehrer and Eilon Solan, \emph{The Maxmin Value of Repeated Games with Incomplete Information on One Side and Tail-Measurable Payoffs}, arXiv:2512.01581 (2025). Locators: Definitions 2.1--2.4 and Theorem 2.5, pp. 6--7; Section 5, p. 18. \url{https://arxiv.org/pdf/2512.01581}
\par\endgroup

\subsection*{Common conventions: Source mathematical conventions}

\textbf{Finite games.} For a finite set $A$, $\Delta(A)$ is its probability
simplex. Players choose independent mixed strategies in Nash equilibrium;
correlation is allowed only when explicitly part of the solution concept.
In a game with payoff functions $u_i$ and strategy profile $x$, an additive
$\varepsilon$ Nash equilibrium satisfies
\[
 u_i(x)\geq u_i(x_i',x_{-i})-\varepsilon
 \quad\text{for every player $i$ and every admissible deviation $x_i'$}.
\]
For numerical approximation, payoffs are normalized as locally specified.
An $\varepsilon$ payoff residual is distinct from distance at most
$\varepsilon$ to the set of exact equilibria. Point convergence, convergence
of empirical play, and convergence in distance to an equilibrium set are
also distinct.

\textbf{Computational representations.} Unless stated otherwise, a complexity
question takes rational data with binary encoding; polynomial time is measured
in total bit length. A fixed-accuracy polynomial algorithm is not necessarily
polynomial in $1/\varepsilon$, and neither is necessarily polynomial in
$\log(1/\varepsilon)$. Succinct games and value, demand, or sampling oracles
must declare their interfaces. Exact real solutions require an explicit
representation; an unspecified list of arbitrary real numbers is not a
finite computational output. A claimed algorithmic barrier is meaningful
only for its specified model and reductions.

\textbf{Dynamic games.} Histories, information, observation, and allowable
randomization are part of the model. A stationary strategy depends only on
the current state; a positional strategy in a deterministic graph game
chooses one action at each controlled vertex. Nash equilibrium at an initial
state and equilibrium after every history are different requirements.
An expectation of a pathwise $\liminf$ payoff, a $\liminf$ of expected averages,
a limiting expected average, and a uniform finite-horizon equilibrium are
different criteria. None is silently substituted for another.

\textbf{Allocation and mechanisms.} A complete allocation assigns every item
exactly once unless otherwise stated. Zero-valued goods, negative values,
capacity constraints, feasibility, lotteries, and copies of goods affect
fairness definitions and are specified locally. Dominant-strategy incentive
compatibility, Bayesian incentive compatibility, universal truthfulness,
and truthfulness in expectation impose different inequalities. Whenever
an objective ratio has zero denominator, its convention is stated or those
instances are excluded.

\textbf{Infinite-dimensional models.} Expectations, integrals, stochastic
equations, and optimization objectives require their local regularity,
integrability, filtrations, and admissible domains. A backward--forward
mean-field system is a boundary-value problem; it is not an initial-value
semiflow without an additional construction. Long-time, large-population,
and vanishing-noise limits require an order of limits and a convergence
topology. If a minimum need not be attained, the objective is its infimum
or supremum and the approximation requirement is stated separately.
% END ECONOMICS STATEMENT
\end{document}
